% Root proof; independent literal audit by D: PASS. Preserved draft in notes.
\subsection{Removing large strong components from minimal candidates}

\begin{lemma}\label{C:extreme-trivial-Rees-block}
Let $C$ be a finite category partitioned into a full strongly
connected source block $S$ and a full complementary category $D$,
with no arrows from $D$ to $S$. Suppose that the minimal ideal
of $S$ has trivial maximal group and exactly one row and one
column at every object. There is an actual functor from $C$ to
the category obtained by replacing $S$ by one trivial object and
each incident right $S$-act by its orbit set. This functor induces
a homotopy equivalence of classifying spaces. The dual statement
holds for a sink block.
\end{lemma}
\begin{proof}
Write $k_{ij}:i\to j$ for the unique minimal-ideal arrow with
these endpoints. The ideal property and uniqueness give
$f k_{hi}=k_{hj}$ and $k_{j\ell}f=k_{i\ell}$ for every
$f:i\to j$ in $S$. In particular the arrows $k_{ij}$ form a
system of compatible zero arrows, and $BS$ is contractible:
the arrows $k_{ia}:i\to a$ give a natural transformation from
the identity functor to the constant functor at a chosen object $a$.

For a finite right $S$-act $F:S^{\mathrm{op}}\to\mathrm{Set}$,
use the category of elements with arrows
$(i,x)\to(j,y)$ represented by $f:i\to j$ with $x=yf$.
Put $e=k_{aa}$ and $T=F(a)e$. Define
\[
 p(i,x)=xk_{ai}\in T,
 \qquad s(t)=(a,t).
\]
Every arrow has the same value under $p$, since
$f k_{ai}=k_{aj}$. Also $ps=\mathrm{id}_T$, and
$k_{ai}$ gives an arrow $s p(i,x)\to(i,x)$ in the element
category. The displayed zero-arrow equations make these arrows
a natural transformation $sp\Rightarrow\mathrm{id}$.
Thus $p$ is a homotopy equivalence onto the discrete set $T$,
which is canonically the orbit set of $F$. The construction is
natural in $F$, including every action of arrows of $D$.

The two-block correspondence description of $BC$ is a homotopy
pushout of the source and target block spaces over their two-sided
category of cross arrows. Applying the preceding natural
equivalence objectwise in $D$ gives a homotopy equivalence of
these cross-arrow spaces to the correspondence for the orbit
category. The source block map $BS\to *$ is an equivalence,
and the target block is unchanged. Homotopy invariance of the
pushout gives the assertion. Composition in the quotient is the
actual induced $D$-action on orbits; there is no path leaving
and then reentering $S$, so no additional composition needs
definition. Opposition proves the dual statement.
\end{proof}

\begin{proposition}\label{C:large-SCC-minimal-candidate21}
Let $C$ be a connected finite category with at most twenty-one
arrows. Assume no object is a retract of another object.
If $C$ has a strongly connected component containing at least
three objects, then $BC$ is not homotopy equivalent to $\CP^2$.
Consequently, a category of type $\CP^2$ with the fewest arrows
and at most twenty-one arrows can have only one- or two-object
strongly connected components.
\end{proposition}
\begin{proof}
A component equal to the whole category is excluded by
Proposition~\ref{C:strong-three-objects-small-cups}.
Let $S$ be a proper component with $t\geq3$ objects.
Because no object is a retract of another, every object identity
lies outside the minimal ideal of $S$. The same adjustment by a
minimal-group inverse as in
Theorem~\ref{C:sixteen-three-SCC} proves this assertion.
Its minimal ideal has at least $t^2$ arrows, so $|S|\geq t^2+t$.
An incident arrow forces an arrow to or from every object of $S$.
The proper complementary category therefore contributes at least
one identity and $t$ cross arrows. For $t\geq4$ this already
requires $t^2+2t+1\geq25$ arrows. Thus $t=3$ and $|S|\geq12$.

If $k$ is the total number of objects, contract $S$ in the
undirected support. Connectivity requires at least $k-3$ support
edges outside $S$, one of which expands to three arrows. With
the remaining identities, this gives
\[
 |C|\geq12+(k-3)+(k-1)=2k+8.
\]
Hence there are at most three objects outside $S$.

If $S$ has both an incoming and an outgoing arrow in the
condensation order, take outside objects $a,b$ with
$a\to S\to b$. They are distinct, since $S$ is a full strongly
connected component. Transitivity supplies three arrows in each
interface and one arrow $a\to b$. Together with $|S|\geq12$
and the two outside identities this costs twenty-one arrows.
Equality is forced. There are no further objects, the outside
corners are trivial and every displayed cross-hom is a singleton.
The object $a$ is initial, so this case is contractible.
We may therefore use opposition to assume that $S$ is a source
block. Write $D=C\setminus S$ and let $T$ be the total number
of cross arrows from $S$ to $D$.

One has $T\geq3$ and $|S|\leq17$. The rational second homology
of $S$ is zero: the categorical Rees Tor formula would otherwise
require total row and column counts at least four, hence at least
sixteen minimal arrows and the three outside identities, giving
$|S|\geq19$. If $H^2(BD;\Q)=0$ as well,
Lemma~\ref{C:two-block-zero-H2-products} makes all degree-two
products of $BC$ zero. Suppose henceforth that $H^2(BD;\Q)\ne0$.

Each connected component of $D$ receives at least one interface
from $S$, hence at least three cross arrows. If $D$ is disconnected,
the component with nonzero second cohomology needs at least five
arrows, by the complete lower-arrow census; another component
needs at least its identity. Thus
$|S|+|D|+T\geq12+6+6=24$, impossible. Therefore $D$ is connected.

If $D$ has three objects, nonzero second cohomology requires at
least eight arrows by Lemma~\ref{C:three-small-H2-zero}; then
$12+8+3>21$. If it has one object with monoid $M$, nonzero
second cohomology requires $|M|\geq5$, whence $|S|\leq13$.
The minimal ideal of $S$ must consequently have trivial group
and total row and column counts both three: the next possible
band size twelve, together with three outside identities, would
already exceed thirteen. It has one row and column per object.
Lemma~\ref{C:extreme-trivial-Rees-block} replaces $S$ by a point
and gives an actual homotopy-equivalent ordinary act cone over
$M$. Its arrow count is at most
\[
 |M|+T+1\leq |C|-|S|+1\leq10.
\]
This contradicts the exact act-cone minimum twenty-two.

It remains that $D$ has two objects. Nonzero second cohomology
requires at least six arrows by
Lemma~\ref{C:two-object-five-H2-zero}. Thus the budget forces
$|S|=12$, $|D|=6$ and $T=3$. As before $S$ has one minimal
row and column per object and trivial minimal group. The same
actual collapse gives a category with three objects and at most
$1+6+3=10$ arrows. Its classifying space is still that of $C$,
contrary to Theorem~\ref{C:CP2-lower16}.
Every case has now been excluded.

Finally a minimal-arrow model cannot contain a retract object:
deleting it preserves the classifying-space type and strictly
decreases the arrow count. The last assertion follows.
\end{proof}
